\centerline{\bf \Large  РЕШЕНИЯ и КРИТЕРИИ}

\centerline{\bf \Large  5 КЛАСС}

\problem{\textbf{1}} Заметим, что 11>9, 111>99, ... ,1111111111>999999999. Соответственно, произведение с единицами больше произведения с девятками. \par

\textbf{Критерии.} \\
\textit{7 баллов} $-$ Произведение с единицами больше, объяснение авторское\\
\textit{1 балл} $-$  Произведение с единицами больше без объяснения\\ 
\textit{0 баллов} $-$ Произведение с девятками больше просто ответ или ответ с объяснением \\

\problem{\textbf{2}} Скорость Арслана $70/7=10$ км/ч, а Данзана $35/7=5$ км/ч. Скорость сближения $10+5=15$ км/ч. Тогда время до встречи $\dfrac{70}{15}$ ч. Тогда расстояние от Оргакина до встречи $\dfrac{70}{15}\times 5 = \dfrac{70}{3} = 23 \dfrac{1}{3}$ км.

Примерное условие Б: Расстояние между Элистой, где живёт Арслан, и Оргакином, где живёт Данзан, равно 70 километров. Данзан пробегает это расстояние за 14 часов. Однажды Арслан и Данзан отправились навстречу друг другу и встретились через $\dfrac{14}{3}$ часа. За сколько часов Арслан пробегает расстояние между Элистой и Оргакином? 

Решение Б: Скорость Данзана $70/14=5$ км/ч. Скорость сближения $70:\dfrac{14}{3}=15$ км/ч. Скорость Арслана $15 - 5 = 10$ км/ч. Арслан пробегает расстояние между Элистой и Оргакином за $70/10=7$ ч\par\par

\textbf{Критерии.} \\
\textit{0-4 баллов} $-$ Решение пункта А\\
\textit{0-1 балла} $-$ Составление условия пункта Б \\
\textit{0-2 баллов} $-$ Решение пункта Б\\
\textit{0 баллов} $-$ Отсутствие решения пункта А (сразу НОЛЬ! баллов)\\
\textit{0 баллов} $-$ Любой ответ без объяснений\\

\problem{\textbf{3}} Заметим, что угловая скорость минутной стрелки больше угловой скорости часовой. Значит, минутная стрелка всегда будет обгонять часовую. В следующий момент после 12:00 минутная стрелка обгонит часовую, и, соответственно, минутная и часовая стрелки не пересекутся в первый час движения в секторе 12-1. Так как минутная стрелка указывает на 12, а часовая - на 1, и за следующий час минутная сделает полный оборот, в то время как часовая продвинется в секторе 1-2. Это означает, что минутная стрелка пересечёт часовую в секторе 1-2 в первый раз, аналогично во второй раз - в секторе 2-3 и так далее. Последнее пересечение произойдёт в секторе 7-8. Пересечение в секторе 7-8 будет, так как минутная за 45 минут дойдёт только до 9. Всего пересечений - 7. Ответ: 7 раз прокричит <<Уралан!>>\par

\centerline{\bf }
\centerline{\bf }
\centerline{\bf }
\centerline{\bf }

\textbf{Критерии.} \\
\textit{7 баллов} $-$ Полностью верное решение\\
\textit{2 балла} $-$ Доказательство того, что минутная и часовая не пересекутся в секторе 12-1\\
\textit{2 балла} $-$ Доказательство того, что минутная и часовая будут пересекаться в секторах 1-2, 2-3 и т.д\\
\textit{2 балла} $-$ Доказательство того, что минутная и часовая пересекутся в секторе 7-8\\
\textit{Снять 1 балл} $-$ Если учитывается <<Уралан!>> в 12:00\\ 
\textit{0-1 балл} $-$ За все остальные попытки решить задачу\\
\textit{0 баллов} $-$ Любой ответ без объяснений\\

\problem{\textbf{4}} Так как прямоугольный лист может быть установлен как вертикально, так и горизонтально, то у нас возможны две ситуации. Пусть $X$ - это количество вертикальных разрезаний, а $Y$ - это количество горизонтальных. Тогда имеем $295	\leqslant 20X 	\leqslant 325$ и $385	\leqslant 25Y 	\leqslant 405$. 

$1!) X=15, Y=16$

$2!) X=16, Y=16$ Или $295	\leqslant 25X 	\leqslant 325$ и $385	\leqslant 20Y 	\leqslant 405$.

$3!) X=12, Y=20$

$4!) X=13, Y=20$

Количество частей на один больше количества разрезаний. А количество прямоугольников - это произведение количества частей после вертикальных разрезаний и количества частей после горизонатальных разрезаний.

$1) 16\cdot17=272$

$2) 17\cdot17=289$

$3) 13\cdot21=273$

$4) 14\cdot21=294$

Ответ: $272, 273, 289, 294$

\textbf{Критерии.} \\
\textit{7 баллов} $-$ Полностью верное решение\\
\textit{По 1 баллу} $-$ За каждое нахождение 1!,2!,3! и 4! \\ 
\textit{1 балл} $-$ Количество частей на 1 больше количества разрезаний\\
\textit{1 балл} $-$ Количество прямоугольников - это произведение количества вертикальных частей на горизонтальных\\
\textit{1 балл} $-$ Указано, что возможны две ситуации в зависимости от изначального расположения листа\\
\textit{1 балл} $-$  За четыре правильных ответа без объяснений\\ 
\textit{0-1 балл} $-$ За все остальные попытки решить задачу\\
\textit{0 баллов} $-$ Меньше четырёх правильных ответов без объяснений\\
\par

\problem{\textbf{5}} Оценка: За 14 часов максимум покрасит в синий $14\cdot7=98$ овец, а нужно покрасить 100. Заметим, что каких бы овец не перекрашивал Очир, количество синих изменяется следующим образом:

+7, +5, +3, +1, -1, -3, -5, -7. 

То есть изменяется на нечётное количество. Значит, за 15 часов изменится также на нечётное количество (как сумма нечётного количества нечётных слагаемых). А изменение должно быть с 0 до 100, то есть чётным. Значит, и за 15 часов тоже не сможет перекрасить в синий все 100 овец. 

Пример: Первые 14-ть часов перекрашиваем 98 овец с красного в синий. За 15-ый час перекрашиваем 6 синих овец в красный, а одну красную в синий. И за 16-ый час перекрашиваем 7 красных в 7 синих. Таким образом, за 16 часов все овцы станут с синей меткой.\par

\textbf{Критерии.} \\
\textit{3 балла} $-$ Правильный пример на 16 часов\\
\textit{1 балл} $-$ Доказательство невозможности за 14 часов\\
\textit{0-3 баллов} $-$ Доказательство невозможности за 15 часов, где 3 балла это правильное доказательство через чётность\\
\textit{0-1 балл} $-$  За все остальные попытки решить задачу\\ 
\textit{0 баллов} $-$ Любой ответ без объяснений\\

